% Einstellungen für das Handout: Layout, Farben, Bausteine und Quellen nach IEEE (deutsch, FHNW-Leitfaden).
% Hier musst du normalerweise nichts ändern. Der Text steht in main.tex, die Quellen in quellen.bib.

\usepackage[T1]{fontenc}
\usepackage[utf8]{inputenc}
\usepackage[nswissgerman]{babel}
\usepackage[autostyle,german=swiss]{csquotes}
\usepackage[default]{sourcesanspro}   % dieselbe Schrift wie auf den Folien (Source Sans)
\usepackage{microtype}
\usepackage[a4paper,left=18mm,right=18mm,top=27mm,bottom=19mm,
            headheight=11mm,headsep=7mm,footskip=9mm]{geometry}
\usepackage{graphicx}
\usepackage{xcolor}
\usepackage{fancyhdr}
\usepackage{titlesec}
\usepackage{enumitem}
\usepackage{tabularx,booktabs,array}
\usepackage{caption}
\usepackage{multicol}
\usepackage[most]{tcolorbox}
\usepackage{fontawesome5}
\usepackage{qrcode}
\usepackage{pgfplots}
\pgfplotsset{compat=1.17}
\usepackage[backend=biber,style=ieee,sorting=none,dashed=false,urldate=long]{biblatex}
\addbibresource{quellen.bib}
\usepackage[hidelinks]{hyperref}

% ---------- Farben: eine Hauptfarbe (Marineblau) und eine Akzentfarbe (Kupfer), wie auf den Folien ----------
% Kupfer steht durchgehend für Gleichstrom und Edison, Marineblau für Wechselstrom und Westinghouse.
\definecolor{navy}{HTML}{10233F}
\definecolor{kupfer}{HTML}{C8641E}
\definecolor{grau}{HTML}{5F6B7A}
\definecolor{hell}{HTML}{EEF1F5}
\definecolor{kupferhell}{HTML}{FBF1E9}
\AtBeginDocument{\color{navy}}

% ---------- Verweise mit noch fehlender Seitenzahl ----------
% \citeoffen{jonnes} ergibt im Entwurf [1, S. ?] in Rot, in der Endfassung [1].
\newif\ifentwurf
\entwurftrue
\newcommand{\citeoffen}[1]{\ifentwurf\cite[{\color{red}S.\,?}]{#1}\else\cite{#1}\fi}
\newcommand{\offen}[1]{\ifentwurf\textcolor{red}{[#1]}\fi}   % sonstige offene Stellen, nur im Entwurf sichtbar

% ---------- Quellenverzeichnis: IEEE mit deutschen Angaben wie im FHNW-Leitfaden ----------
\stdpunctuation                                                    % «Titel», Zeitschrift (Komma nach dem Schlusszeichen)
\DefineBibliographyStrings{nswissgerman}{jourvol = {Bd\adddot}}   % Bd. statt Jg.
\DeclareFieldFormat{url}{\mkbibbrackets{Online}\adddot\space Verfügbar unter\addcolon\space\url{#1}}
\DeclareFieldFormat{urldate}{\mkbibbrackets{Zugegriffen\addcolon\space#1}}
\renewcommand*{\bibfont}{\footnotesize\raggedright}
\setlength{\bibitemsep}{1.5pt}
% Webseiten: Urheber, «Titel», Organisation, Datum. [Online]. Verfügbar unter: Adresse [Zugegriffen: Datum].
\DeclareBibliographyDriver{online}{%
  \usebibmacro{bibindex}%
  \usebibmacro{begentry}%
  \usebibmacro{author/editor+others/translator+others}%
  \setunit{\labelnamepunct}\newblock
  \usebibmacro{title}%
  \newunit\newblock
  \printlist{organization}%
  \newunit\newblock
  \usebibmacro{date}%
  \setunit{\adddot\addspace}\newblock
  \usebibmacro{url+urldate}%
  \usebibmacro{finentry}}

% ---------- Seitenlayout ----------
\pagestyle{fancy}
\fancyhf{}
\fancyhead[L]{\includegraphics[height=9mm]{bilder/fhnw_logo.jpg}}
\fancyhead[R]{\footnotesize\color{grau}Modul «Weltmacht USA» \textperiodcentered{} Handout}
\fancyfoot[L]{\footnotesize\color{grau}FHNW, School of Engineering, Windisch \textperiodcentered{} Dozent: Stephan Schwarz}
\fancyfoot[R]{\footnotesize\color{grau}\thepage}
\renewcommand{\headrulewidth}{0.4pt}
\renewcommand{\footrulewidth}{0pt}
\renewcommand{\headrule}{{\color{navy!30}\hrule width\headwidth height\headrulewidth}}

\setlength{\parindent}{0pt}
\setlength{\emergencystretch}{1.5em}
\setlength{\parskip}{4pt plus 1pt}
\setcounter{secnumdepth}{0}
% Überschriften mit kurzem Kupferstrich
\titleformat{\section}{\large\bfseries\color{navy}}{}{0pt}{}[{\vspace{-3pt}\textcolor{kupfer}{\rule{11mm}{1.3pt}}}]
\titlespacing*{\section}{0pt}{11pt plus 2pt minus 2pt}{3pt}
\setlist{nosep,leftmargin=*}
\renewcommand{\arraystretch}{1.25}
\newcolumntype{L}[1]{>{\raggedright\arraybackslash}p{#1}}
\newcolumntype{Y}{>{\raggedright\arraybackslash}X}
\captionsetup{labelsep=colon,skip=4pt,font={footnotesize,color=grau},labelfont={bf,color=navy},justification=raggedright,singlelinecheck=false}
\captionsetup[table]{name=Tab.,position=bottom}
\captionsetup[figure]{name=Abb.}

% ---------- Bausteine ----------
% Leitfrage (dunkel) und Kurzantwort (hell)
\newtcolorbox{leitfrage}{enhanced,colback=navy,colframe=navy,coltext=white,arc=1.5mm,boxrule=0pt,
  left=4mm,right=4mm,top=2.6mm,bottom=2.6mm,fontupper=\normalsize,
  colbacklower=hell,collower=navy,fontlower=\normalsize,lowerbox=visible,segmentation style={draw=none},
  bicolor,boxsep=0pt,middle=2.6mm}
\newcommand{\etikett}[1]{{\footnotesize\bfseries\color{kupfer}\MakeUppercase{#1}}\par\vspace{1pt}}
% Karte mit farbigem Balken oben: \begin{karte}{Farbe}{Titel} ... \end{karte}
\newtcolorbox{karte}[2]{enhanced,colback=white,colframe=#1,boxrule=0.4pt,toprule=2.2pt,arc=1.2mm,
  left=3mm,right=3mm,top=2mm,bottom=2mm,boxsep=0pt,
  title={#2},fonttitle=\bfseries,coltitle=navy,attach title to upper={\par\vspace{2pt}}}
% Kachel mit grosser Zahl: \kennzahl{Farbe}{Zahl}{Erklärung}
\newcommand{\kennzahl}[3]{\begin{tcolorbox}[enhanced,colback=hell,colframe=hell,arc=1.2mm,boxrule=0pt,
  left=3mm,right=3mm,top=2mm,bottom=2mm,boxsep=0pt]
  {\fontsize{24}{26}\selectfont\bfseries\color{#1}#2}\par\vspace{1pt}{\footnotesize\raggedright #3\par}\end{tcolorbox}}
% Nummerierter Punkt: \punkt{1}{Titel}{Text}
\newcommand{\punkt}[3]{\begin{tcolorbox}[enhanced,colback=white,colframe=navy!18,boxrule=0.4pt,arc=1.2mm,
  left=3mm,right=3mm,top=2mm,bottom=2mm,boxsep=0pt]
  {\fontsize{20}{20}\selectfont\bfseries\color{kupfer}#1}\hspace{2mm}{\bfseries #2}\par\vspace{2pt}{\small\raggedright #3\par}\end{tcolorbox}}
% Zeile mit Symbol: \symbolzeile{\faBolt}{Titel}{Text}
\newcommand{\symbolzeile}[3]{\par\noindent\begin{minipage}[t]{6mm}\centering\color{kupfer}#1\end{minipage}%
  \begin{minipage}[t]{\dimexpr\linewidth-6mm}\raggedright\textbf{#2} #3\end{minipage}\par\vspace{3pt}}
% QR-Code mit Beschriftung: \qr{Adresse}{Titel}{Beschreibung}
\newcommand{\qr}[3]{\begin{minipage}[t]{15mm}\vspace{0pt}\qrcode[height=14mm]{#1}\end{minipage}\hspace{2.5mm}%
  \begin{minipage}[t]{\dimexpr\linewidth-17.5mm}\vspace{0pt}\raggedright\small\textbf{#2}\par #3\end{minipage}}

% Ereignis auf dem Zeitstrahl: \ereignis{Jahr}{oben|unten}{Titel}{Text}; 1882 liegt bei x = 15 mm, ein Jahr = 10,286 mm
\newcommand{\zeitx}[1]{15+(#1-1882)*10.286}
\newcommand{\ereignis}[4]{%
  \ifstrequal{#2}{oben}
    {\draw[navy!35,line width=0.5pt] ({\zeitx{#1}},1.4) -- ++(0,3.6);
     \node[anchor=south,align=center,text width=29mm,inner sep=0.8pt,font=\footnotesize]
       at ({\zeitx{#1}},5.2) {{\bfseries\color{kupfer}#1}\\[-1.5pt]\textbf{#3}\\[-1.5pt]{\scriptsize\color{grau}#4}};}
    {\draw[navy!35,line width=0.5pt] ({\zeitx{#1}},-1.4) -- ++(0,-3.6);
     \node[anchor=north,align=center,text width=29mm,inner sep=0.8pt,font=\footnotesize]
       at ({\zeitx{#1}},-5.2) {{\bfseries\color{kupfer}#1}\\[-1.5pt]\textbf{#3}\\[-1.5pt]{\scriptsize\color{grau}#4}};}
  \fill[white] ({\zeitx{#1}},0) circle (1.9);
  \fill[navy] ({\zeitx{#1}},0) circle (1.25);}
% Begriff im Glossar
\newcommand{\begriff}[2]{\par\noindent\textbf{#1}\quad #2\par\vspace{3pt}}

% ---------- Grafiken: Zahlen in der Textschrift ----------
\pgfplotsset{/pgf/number format/assume math mode=true}
